\documentclass[12pt]{article}
\usepackage[utf8]{inputenc}
\usepackage[T1]{fontenc}
\usepackage[margin=1in]{geometry}
\usepackage{setspace}
\usepackage{xcolor}
\usepackage[colorlinks=true, urlcolor=blue!60!black, linkcolor=blue!60!black]{hyperref}
\setlength{\parindent}{0em}
\setlength{\parskip}{0.6em}
\renewcommand{\baselinestretch}{1.0}
\begin{document}
\begin{center}
{\LARGE \textit{Mathews v. Eldridge}, 424 U.S. 319 (1976)\par}
\vspace{0.5em}
{\large Case Brief --- Two-Page Edition\par}
\end{center}

\section{Memory Jogger}
Disability benefits may be terminated before an evidentiary hearing; procedural due process is governed by a three-factor balancing test (\href{https://supreme.justia.com/cases/federal/us/424/319/}{424 U.S. 319}).

\section{Detailed Case Facts}
Social Security disability benefits are not means-tested. Before terminating them, a state agency gives written notice of a tentative decision, a summary of the medical evidence, access to the file, and a chance for written rebuttal; afterward the claimant may get an evidentiary hearing and retroactive benefits. George Eldridge's benefits were terminated in 1972 after the agency found his disability had ceased.

\section{Procedural History}
Relying on \textit{Goldberg v.\ Kelly}, the district court required a pre-termination hearing; the Fourth Circuit affirmed; the Supreme Court reversed.

\section{Judicial Votes}
6--2. Powell wrote for Burger, Stewart, White, Blackmun, and Rehnquist; Brennan dissented, joined by Marshall; Stevens did not participate.

\section{Holding}
The district court had jurisdiction. Due process requires weighing the private interest, the risk of error and value of added safeguards, and the government's interest; the SSA's written pre-termination procedures plus a post-termination hearing suffice.

\section{Analysis of Opinions}
Powell found the private interest less weighty than in \textit{Goldberg} because disability benefits are not need-based and back pay is available. Medical determinations rest on standard physicians' reports, not credibility, so written procedures reduce error; rules are shaped by the generality of cases. Added hearings would cost money that may come from the deserving. Brennan called the deprivation's limits speculative, noting that Eldridge's family lost its home.

\section{Examples: Future Applications}
\textbf{Same-Side:} reducing a veteran's disability rating on medical evidence after written notice, with a later hearing. \textbf{Opposite-Side:} cutting off subsistence welfare over a credibility dispute requires a prior hearing (\textit{Goldberg}). \textbf{Fence-Sitter:} algorithmic cuts to in-home Medicaid nursing hours.

\section{Critique}
Jerry Mashaw faulted \textit{Mathews} as cost-benefit accounting that ignores dignitary values; originalists call it ahistorical.

\section{Key Quotations}
\begin{quote}
``More precisely, our prior decisions indicate that identification of the specific dictates of due process generally requires consideration of three distinct factors: first, the private interest that will be affected by the official action; second, the risk of an erroneous deprivation of such interest through the procedures used, and the probable value, if any, of additional or substitute procedural safeguards; and, finally, the Government's interest, including the function involved and the fiscal and administrative burdens that the additional or substitute procedural requirement would entail.''
\end{quote}
\noindent --- Justice Powell, Opinion of the Court, 424 U.S. at 334--35.

\end{document}
